\section*{Research Experience}
My established and exploratory research topics, each with a representative mix of publications, projects, and presentations.
\listsnote
\researchplannote

\subsection*{Trajectory Representation Learning}
Trajectory representation learning is one of my established research topics. It aims to learn representations of trajectories (of vehicle, human, or vessel) that are universal across various downstream tasks. This also extends to building a multi-task or foundational model for trajectories that can perform various tasks at once.

\pubentry{IEEE TKDE, 2025, Impact Factor 11.6}{\href{https://ieeexplore.ieee.org/document/10818577}{Paper} \,/\, \href{https://github.com/Logan-Lin/UniTE}{Code}}
{UniTE: A Survey and Unified Pipeline for Pre-training Spatiotemporal Trajectory Embeddings}
{\me, Zeyu Zhou, Yicheng Liu, Haochen Lv, Haomin Wen, Tianyi Li, Yushuai Li, Christian S. Jensen, Shengnan Guo, Youfang Lin, Huaiyu Wan}
{Surveys existing methods for learning reusable representations of movement trajectories and brings them together into a single modular framework with shared code, so that new methods can be built, compared, and evaluated on the same basis.}

\pubentry{IEEE TKDE, 2025, Impact Factor 11.6}{\href{https://ieeexplore.ieee.org/document/11004614}{Paper} \,/\, \href{https://github.com/Logan-Lin/UVTM}{Code}}
{UVTM: Universal Vehicle Trajectory Modeling with ST Feature Domain Generation}
{\me, Jilin Hu, Shengnan Guo, Bin Yang, Christian S. Jensen, Youfang Lin, Huaiyu Wan}
{A single model for vehicle GPS trajectories that handles many tasks such as travel time estimation, trajectory recovery, and trajectory prediction, instead of maintaining a separate model for each. It stays accurate even when trajectories are sparse or only part of their features are available, by learning to rebuild dense and complete trajectories from incomplete ones.}

\pubentry{NeurIPS, 2025, CORE Rank A*}{\href{https://neurips.cc/virtual/2025/oral/117509}{Oral} \,/\, \href{https://neurips.cc/virtual/2025/poster/117508}{Poster}}
{TransferTraj: A Vehicle Trajectory Learning Model for Region and Task Transferability}
{Tonglong Wei*, \me*, Zeyu Zhou, Haomin Wen, Jilin Hu, Shengnan Guo, Youfang Lin, Gao Cong, Huaiyu Wan}
{Learns from vehicle GPS trajectories in a way that transfers across different geographic regions and different prediction tasks without retraining, removing the need to keep separate specialized models. Handles each task by treating it as recovering hidden parts of a trajectory, so one model serves many tasks even with limited data.}

\pubentry{IEEE TKDE, 2023, Impact Factor 8.9}{\href{https://ieeexplore.ieee.org/abstract/document/10375102}{Paper} \,/\, \href{https://github.com/Logan-Lin/MMTEC}{Code}}
{Pre-training General Trajectory Embeddings with Maximum Multi-view Entropy Coding}
{\me, Huaiyu Wan, Shengnan Guo, Jilin Hu, Christian S. Jensen, Youfang Lin}
{Learns general-purpose representations of movement trajectories from unlabeled data that capture both travel behavior and spatial and temporal patterns. The learned representations avoid task-specific bias so they transfer well across many downstream tasks.}

\pubentry{AAAI, 2021, CORE Rank A*}{\href{https://ojs.aaai.org/index.php/AAAI/article/view/16548}{Paper} \,/\, \href{https://github.com/Logan-Lin/CTLE}{Code}}
{Pre-training Context and Time Aware Location Embeddings from Spatial-Temporal Trajectories for User Next Location Prediction}
{\me, Huaiyu Wan, Shengnan Guo, Youfang Lin}
{Pre-trains location representations from movement trajectories that capture how the meaning of a place changes with its surrounding context and the time of visit, leading to more accurate prediction of a user's next location.}

\subsection*{Spatiotemporal Data Mining}
Spatiotemporal data mining is a broader topic that is inclusive of the above. It covers research on extracting and utilizing patterns from large-scale data with spatial information that is dynamic over time, usually sourced from transportation scenarios.

\pubentry{IJCAI, 2026, CORE Rank A*}{\href{https://arxiv.org/abs/2603.29837}{Preprint} \,/\, \href{https://github.com/Logan-Lin/DiSGMM-code}{Code}}
{DiSGMM: A Method for Time-varying Microscopic Weight Completion on Road Networks}
{\me, Jilin Hu, Shengnan Guo, Christian S. Jensen, Youfang Lin, Huaiyu Wan}
{Fills in missing fine-grained, time-varying traffic conditions on road networks, such as travel speeds on individual road segments during specific time periods, when observations are sparse both across segments and within each segment. Estimates the full range of likely conditions for each segment and time rather than a single value.}

\pubentry{NeurIPS, 2025, CORE Rank A*}{\href{https://neurips.cc/virtual/2025/poster/117945}{Poster} \,/\, \href{https://arxiv.org/abs/2410.14281}{Preprint}}
{PLMTrajRec: A Scalable and Generalizable Trajectory Recovery Method with Pre-trained Language Models}
{Tonglong Wei*, \me*, Youfang Lin, Shengnan Guo, Jilin Hu, Haitao Yuan, Gao Cong, Huaiyu Wan}
{Recovers the missing points in sparse movement trajectories to restore detailed paths, while needing only a small amount of dense training data and generalizing across trajectories recorded at different sampling rates.}

\pubentry{WWW, 2025, CORE Rank A*}{\href{https://openreview.net/forum?id=KmMSQS6tFn}{Paper} \,/\, \href{https://github.com/decisionintelligence/Path-LLM}{Code}}
{Path-LLM: A Multi-Modal Path Representation Learning by Aligning and Fusing with Large Language Models}
{Yongfu Wei*, \me*, Hongfan Gao, Ronghui Xu, Sean Bin Yang, Jilin Hu}
{Learns representations of paths in a road network by combining the network structure with text that describes physical and regional context, which earlier methods left out. The combined view improves accuracy on tasks such as ranking paths and estimating travel time, including settings with little or no labeled data.}

\pubentry{AAAI, 2025, CORE Rank A*}{\href{https://ojs.aaai.org/index.php/AAAI/article/view/33350}{Paper} \,/\, \href{https://arxiv.org/abs/2408.12809}{Preprint}}
{DutyTTE: Deciphering Uncertainty in Origin-Destination Travel Time Estimation}
{Xiaowei Mao*, \me*, Shengnan Guo, Yubin Chen, Xingyu Xian, Haomin Wen, Qisen Xu, Youfang Lin, Huaiyu Wan}
{Estimates the travel time between an origin and a destination and measures how uncertain each prediction is, by first inferring the likely route the trip will follow and then producing reliable confidence intervals that reflect how individual road segments contribute to the overall uncertainty under varying conditions.}

\pubentry{SIGMOD, 2024, CORE Rank A*}{\href{https://dl.acm.org/doi/10.1145/3617337}{Paper} \,/\, \href{https://github.com/Logan-Lin/DOT}{Code}}
{Origin-Destination Travel Time Oracle for Map-based Services}
{\me, Huaiyu Wan, Jilin Hu, Shengnan Guo, Bin Yang, Christian S. Jensen, Youfang Lin}
{Estimates the travel time between an origin and a destination at a given departure time by learning from many past trips that connect the same pair of locations. It first infers a likely route between the two points and then predicts the travel time along it, improving accuracy for map-based navigation services.}

\subsection*{AI for Materials Science}
AI for materials science is a new direction I am exploring. The promise is that AI enables discovery of new materials with tailored properties, a difficult process in traditional materials design.

\pubentry{Advanced Materials, 2026, Impact Factor 29.1}{\href{https://doi.org/10.1002/adma.202522493}{Paper} \,/\, \href{https://arxiv.org/abs/2509.13916}{Preprint} \,/\, \href{https://github.com/Logan-Lin/AMDEN-code}{Code}}
{Inverse Design of Amorphous Materials with Targeted Properties}
{Jonas A. Finkler*, \me*, Tao Du, Jilin Hu, Morten M. Smedskjaer}
{Generates atomic structures of disordered materials such as glasses that match desired target properties, and refines them into stable low-energy configurations. Also introduces new datasets of amorphous materials to support this kind of design.}

\pubentry{NeurIPS Workshop, 2025}{\href{https://openreview.net/forum?id=cEgjPFdLvl}{Paper}}
{AMDEN: Amorphous Materials DEnoising Network}
{Jonas A. Finkler*, \me*, Morten M. Smedskjaer, Tao Du, Jilin Hu}{}

\pubentry{PNCS17, Oral Presentation, 2025}{\href{https://github.com/Logan-Lin/amden-presentation/releases/latest/download/slide.pdf}{Slides}}
{AMDEN: Amorphous Materials DEnoising Network}
{\me*, Jonas A. Finkler*, Tao Du, Morten M. Smedskjaer, Jilin Hu}{}

\pubentry{AAAI, 2026, CORE Rank A*}{\href{https://ojs.aaai.org/index.php/AAAI/article/view/37139}{Paper} \,/\, \href{https://arxiv.org/abs/2511.12489}{Preprint}}
{SculptDrug: A Spatial Condition-Aware Bayesian Flow Model for Structure-based Drug Design}
{Qingsong Zhong, Haomin Yu, \me, Wangmeng Shen, Long Zeng, Jilin Hu}
{Generates drug molecules that fit a target protein's three-dimensional structure, keeping the molecules shaped to sit within the protein's surface and consistent with both its overall form and fine details. This produces more accurate candidate molecules for structure-based drug discovery.}

\pubentry{Villum Foundation, Project}{}
{Inverse Design of Materials Using Diffusion Probabilistic Models}
{}
{Diffusion probabilistic models for understanding chemistry-structure-property relationships and inverse design of new materials.}
